% GENERATED FILE. Edit canonical lessons, then run npm run generate:books.
% canonical-source-hash: fnv1a64:858a4f857e7b5095
% canonical-lessons: MR-C243-thikan, MR-C243-dekhava, MR-C243-sahal, MR-C243-tava, MR-C243-kadhai

\chapter{Place, View, Outing, Griddle, Wok}
\label{ch:mr-place-view-outing-griddle-wok}
\hlchaptermodality{243}

\begin{chapteropening}
\textbf{By the end of this chapter:} \emph{I can say a place, a view, an outing, a griddle and a wok.}

Complete the last lesson of chapter 243: I can say a place, a view, an outing, a griddle and a wok.
\end{chapteropening}

\section[\texorpdfstring{\mr{ठिकाण} (ṭhikāṇ)}{ठिकाण (ṭhikāṇ)}]{\mr{ठिकाण} (ṭhikāṇ) — a place, a spot}

\label{lesson:MR-C243-thikan}



\begin{quote}
Before the new one: say the Marathi for a fine thread, then the Marathi for a ladder.
\end{quote}

\subsection*{You'll want to know: \mr{ठिकाण}}
\textbf{\mr{ठिकाण}} — \emph{ṭhikāṇ} — \textquotedblleft{}a place, a spot\textquotedblright{}.

\begin{grammarlens}[title={What you've built}]
One more word for this chapter.
\end{grammarlens}

\subsection*{Your turn}
\begin{itemize}
\raggedright
  \item \emph{Say it:} \emph{ṭhikāṇ}
  \item \emph{Say it:} \emph{ṭhikāṇ}, once more
  \item \emph{Say it:} say \emph{ṭhikāṇ}
  \item \emph{Recall:} say the Marathi for a fine thread, then the Marathi for a ladder, then say \emph{ṭhikāṇ} again
\end{itemize}

\begin{tcolorbox}[breakable,colback=teal!4,colframe=teal!35!black,title={Before you move on}]
What does \mr{ठिकाण} mean? (\textquotedblleft{}A place, a spot\textquotedblright{}.) Say it once more.
\end{tcolorbox}
\section[\texorpdfstring{\mr{देखावा} (dekhāvā)}{देखावा (dekhāvā)}]{\mr{देखावा} (dekhāvā) — a view, a scene}

\label{lesson:MR-C243-dekhava}



\begin{quote}
Before the new one: say the Marathi for a ladder, then the Marathi for a place.
\end{quote}

\subsection*{You'll want to know: \mr{देखावा}}
\textbf{\mr{देखावा}} — \emph{dekhāvā} — \textquotedblleft{}a view, a scene\textquotedblright{}.

\begin{grammarlens}[title={What you've built}]
One more word for this chapter.
\end{grammarlens}

\subsection*{Your turn}
\begin{itemize}
\raggedright
  \item \emph{Say it:} \emph{dekhāvā}
  \item \emph{Say it:} \emph{dekhāvā}, once more
  \item \emph{Say it:} say \emph{dekhāvā}
  \item \emph{Recall:} say the Marathi for a ladder, then the Marathi for a place, then say \emph{dekhāvā} again
\end{itemize}

\begin{tcolorbox}[breakable,colback=teal!4,colframe=teal!35!black,title={Before you move on}]
What does \mr{देखावा} mean? (\textquotedblleft{}A view, a scene\textquotedblright{}.) Say it once more.
\end{tcolorbox}
\section[\texorpdfstring{\mr{सहल} (sahal)}{सहल (sahal)}]{\mr{सहल} (sahal) — an outing, a trip}

\label{lesson:MR-C243-sahal}



\begin{quote}
Before the new one: say the Marathi for a place, then the Marathi for a view.
\end{quote}

\subsection*{You'll want to know: \mr{सहल}}
\textbf{\mr{सहल}} — \emph{sahal} — \textquotedblleft{}an outing, a trip\textquotedblright{}.

\begin{grammarlens}[title={What you've built}]
One more word for this chapter.
\end{grammarlens}

\subsection*{Your turn}
\begin{itemize}
\raggedright
  \item \emph{Say it:} \emph{sahal}
  \item \emph{Say it:} \emph{sahal}, once more
  \item \emph{Say it:} say \emph{sahal}
  \item \emph{Recall:} say the Marathi for a place, then the Marathi for a view, then say \emph{sahal} again
\end{itemize}

\begin{tcolorbox}[breakable,colback=teal!4,colframe=teal!35!black,title={Before you move on}]
What does \mr{सहल} mean? (\textquotedblleft{}An outing, a trip\textquotedblright{}.) Say it once more.
\end{tcolorbox}
\section[\texorpdfstring{\mr{तवा} (tavā)}{तवा (tavā)}]{\mr{तवा} (tavā) — a griddle}

\label{lesson:MR-C243-tava}



\begin{quote}
Before the new one: say the Marathi for a view, then the Marathi for an outing.
\end{quote}

\subsection*{You'll want to know: \mr{तवा}}
\textbf{\mr{तवा}} — \emph{tavā} — \textquotedblleft{}a griddle\textquotedblright{}.

\begin{grammarlens}[title={What you've built}]
One more word for this chapter.
\end{grammarlens}

\subsection*{Your turn}
\begin{itemize}
\raggedright
  \item \emph{Say it:} \emph{tavā}
  \item \emph{Say it:} \emph{tavā}, once more
  \item \emph{Say it:} say \emph{tavā}
  \item \emph{Recall:} say the Marathi for a view, then the Marathi for an outing, then say \emph{tavā} again
\end{itemize}

\begin{tcolorbox}[breakable,colback=teal!4,colframe=teal!35!black,title={Before you move on}]
What does \mr{तवा} mean? (\textquotedblleft{}A griddle\textquotedblright{}.) Say it once more.
\end{tcolorbox}
\section[\texorpdfstring{\mr{कढई} (kaḍhaī)}{कढई (kaḍhaī)}]{\mr{कढई} (kaḍhaī) — a wok}

\label{lesson:MR-C243-kadhai}



\begin{quote}
Before the new one: say the Marathi for an outing, then the Marathi for a griddle.
\end{quote}

\subsection*{You'll want to know: \mr{कढई}}
\textbf{\mr{कढई}} — \emph{kaḍhaī} — \textquotedblleft{}a wok\textquotedblright{}.

\begin{grammarlens}[title={What you've built}]
One more word for this chapter.
\end{grammarlens}

\subsection*{Your turn}
\begin{itemize}
\raggedright
  \item \emph{Say it:} \emph{kaḍhaī}
  \item \emph{Say it:} \emph{kaḍhaī}, once more
  \item \emph{Say it:} say \emph{kaḍhaī}
  \item \emph{Recall:} say the Marathi for an outing, then the Marathi for a griddle, then say \emph{kaḍhaī} again
\end{itemize}

\begin{tcolorbox}[breakable,colback=teal!4,colframe=teal!35!black,title={Before you move on}]
What does \mr{कढई} mean? (\textquotedblleft{}A wok\textquotedblright{}.) Say it once more.
\end{tcolorbox}
